\section{波形信道}

\begin{remarkBox}{本章重点}
本章把第六章的信道容量从离散符号推广到连续取值和连续时间波形。学习主线是：先建立时间离散连续信道模型，再用约束容量处理连续输入的功率限制；随后分析加性噪声信道和加性高斯噪声信道，掌握 AWGN 容量公式；最后把限带波形信道等价成 $N=2TW$ 个自由度的并联信道，并理解功率利用率、频谱利用率和 Shannon 极限之间的关系。
\end{remarkBox}

前面几章主要研究离散信道和编码定理。本章面对的信号更接近实际通信系统：输入、噪声和输出可以是连续取值的随机变量，也可以是随时间变化的波形。处理这类信道时，核心方法不是直接在连续时间上硬算，而是把波形信道离散化成许多独立的自由度，再沿用互信息最大化和容量的思想。

需要把握的一条线索是：连续信道若不限制输入功率或幅度，容量往往会无限大。因此本章所有容量公式都离不开约束条件。对 AWGN 信道来说，平均功率约束和高斯噪声共同给出著名的 Shannon 公式，这也是通信系统设计中最常用的理论上限。

\subsection{时间离散连续信道}

\subsubsection{信道模型}

时间离散连续信道的输入和输出按离散时刻出现，但每个符号的取值可以是连续的。它可以来自对连续时间信道的抽样，也可以来自某种正交变换后的系数信道。和离散信道类似，先研究单符号信道，再推广到多维矢量信道。

设信道输入和输出分别为长为 $N$ 的随机矢量
\[
X^N=(X_1,X_2,\dots,X_N),
\qquad
Y^N=(Y_1,Y_2,\dots,Y_N),
\]
对应取值为
\[
x=(x_1,x_2,\dots,x_N),
\qquad
y=(y_1,y_2,\dots,y_N).
\]
一般时间离散连续信道写成
\[
\{X^N,p(y\mid x),Y^N\},
\]
其中
\[
p(y\mid x)=p(y_1,y_2,\dots,y_N\mid x_1,x_2,\dots,x_N)
\]
是信道转移概率密度。

\begin{definitionBox}{时间离散无记忆连续信道}
若信道转移概率密度满足
\[
p(y\mid x)=\prod_{n=1}^{N}p(y_n\mid x_n),
\]
则称该信道为时间离散无记忆连续信道。

若进一步对任意时刻 $m,n$ 都有
\[
p(y_n\mid x_n)=p(y_m\mid x_m),
\]
则称为平稳或恒参无记忆连续信道。此时单符号模型可写成
\[
\{X,p(y\mid x),Y\}.
\]
\end{definitionBox}

无记忆分解的含义和离散信道相同：每个时刻的输出只依赖当前输入，不依赖过去输入输出。平稳性则说明同一个单符号转移概率密度在每个时刻都适用。除非题目特别说明，本章后面研究的连续信道都默认是平稳的。

\subsubsection{连续信道容量的约束}

离散信道容量只需要在输入概率分布上最大化；连续信道不同，若不限制输入幅度或功率，互信息最大值可能无限增大。因此连续信道容量通常定义为满足某种输入约束时的最大平均互信息。

\begin{definitionBox}{约束容量}
设 $f(x)$ 是与输入有关的非负代价函数，$\beta$ 是约束量。单符号时间离散连续信道的约束容量定义为
\[
C(\beta)=\max_{p(x):\,E[f(X)]\le\beta} I(X;Y).
\]
通常容量随 $\beta$ 增大而增大，因此求最大值时常在
\[
E[f(X)]=\beta
\]
处达到。
\end{definitionBox}

最常见的约束是平均能量约束。取
\[
f(x)=x^2,
\qquad
\beta=E,
\]
就得到
\[
E[X^2]\le E.
\]
在这种约束下，平稳无记忆连续信道的容量写成
\[
C=\max_{p(x):\,E[X^2]\le E} I(X;Y).
\]

若输入信源无记忆，多维信道的每自由度容量也可归结到单符号容量：
\[
\frac1N\max_{p(x^N)}I(X^N;Y^N)=\max_{p(x)}I(X;Y).
\]
这说明对平稳无记忆连续信道，关键仍是找出单符号信道在约束下的最优输入分布。

\subsection{加性噪声信道与容量}

\subsubsection{加性噪声模型}

加性噪声信道是本章最重要的连续信道模型。它假设信道输出等于输入加上噪声，噪声独立于输入。很多实际系统经过等效建模后都可以先近似成这个形式。

设输入为 $X$，噪声为 $Z$，输出为
\[
Y=X+Z,
\]
其中 $Z$ 独立于 $X$，概率密度为 $p_Z(z)$。由于给定 $X=x$ 时，$Y=x+Z$，所以信道转移概率密度为
\[
p(y\mid x)=p_Z(y-x).
\]

\begin{theoremBox}{加性噪声信道的互信息}
对加性噪声信道 $Y=X+Z$，若噪声 $Z$ 独立于输入 $X$，则
\[
h(Y\mid X)=h(Z),
\]
从而
\[
I(X;Y)=h(Y)-h(Z).
\]
因此在噪声分布固定时，求容量等价于在输入约束下最大化输出差熵 $h(Y)$：
\[
C=\max_{p(x)}\{h(Y)-h(Z)\}.
\]
\end{theoremBox}

这个公式的直觉很清楚：噪声的不确定性 $h(Z)$ 是信道本身带来的代价，输入分布能做的是让输出 $Y$ 在约束范围内尽量分散，从而让 $h(Y)$ 尽量大。

\begin{example}
\label{ex:ch8-amplitude-limited-uniform-noise}
\textbf{幅度受限的均匀噪声信道}：设噪声
\[
Z\sim U[-1,1],
\]
输入满足幅度约束
\[
X\in[-1,1].
\]
求输入与输出平均互信息的最大值。

\noindent \textbf{解题流程}：\\[2pt]
步骤一：先由 $Y=X+Z$ 判断输出的取值范围。\\
步骤二：噪声熵固定，最大化 $I(X;Y)$ 等价于最大化 $h(Y)$。\\
步骤三：在幅度受限范围内用均匀分布最大熵结论计算差熵差。

\textbf{解}：因为 $X\in[-1,1]$，$Z\in[-1,1]$，所以
\[
Y\in[-2,2].
\]
在给定取值区间时，均匀分布具有最大差熵。因此
\[
\max I(X;Y)=\max h(Y)-h(Z)
=\log_2 4-\log_2 2=1.
\]
所以最大平均互信息为
\[
1\text{ bit/自由度}.
\]
\end{example}

\subsubsection{离散时间 AWGN 信道容量}

若加性噪声 $Z$ 服从高斯分布，就得到加性高斯噪声信道。高斯噪声特别重要，因为在相同功率下，高斯噪声使信道容量最小，是最难抵抗的加性噪声。

设
\[
Z\sim N(0,\sigma_z^2),
\qquad
E[X^2]\le E.
\]
当输入 $X$ 也是零均值高斯随机变量，且方差为 $E$ 时，输出
\[
Y=X+Z
\]
也是高斯随机变量，方差为 $E+\sigma_z^2$。由高斯分布最大差熵性质可得容量。

\begin{formulaBox}{离散时间 AWGN 信道容量}
对离散时间平稳无记忆加性高斯噪声信道，若噪声方差为 $\sigma_z^2$，输入平均能量约束为 $E[X^2]\le E$，则每自由度容量为
\[
C=\frac12\log_2\left(1+\frac{E}{\sigma_z^2}\right)
\quad\text{bit/自由度}.
\]
达到容量时，输入为满足功率约束的高斯分布。
\end{formulaBox}

这里的 $1/2$ 来自一维实高斯随机变量的差熵公式。容量只与输入信噪比 $E/\sigma_z^2$ 有关。若不限制 $E$，容量会随信噪比增大而无界增长。

\subsection{并联高斯信道与注水分配}

实际信道常常可以通过频域分解或正交展开变成多个并联子信道。若各子信道噪声方差不同，总功率应该优先分给噪声小的子信道，这就是注水原理。

设共有 $N$ 个独立高斯子信道：
\[
Y_i=X_i+Z_i,
\qquad
Z_i\sim N(0,\sigma_i^2),
\qquad i=1,2,\dots,N.
\]
输入能量分配为 $E_i=E[X_i^2]$，满足
\[
\sum_{i=1}^{N}E_i\le E.
\]
并联信道容量为
\[
C=\sum_{i=1}^{N}\frac12\log_2\left(1+\frac{E_i}{\sigma_i^2}\right),
\]
其中 $E_i$ 由总能量约束下的最大化决定。

\begin{theoremBox}{并联高斯信道的注水原理}
达到容量时，各子信道输入统计独立且服从高斯分布。能量分配满足
\[
E_i=\begin{cases}
B-\sigma_i^2, & B>\sigma_i^2,\\
0, & B\le\sigma_i^2,
\end{cases}
\]
并由
\[
\sum_{i:B>\sigma_i^2}(B-\sigma_i^2)=E
\]
确定水位 $B$。
\end{theoremBox}

可以把 $\sigma_i^2$ 看成第 $i$ 个池底高度，把总能量 $E$ 看成水量。水位 $B$ 高于池底的部分就是分配给该子信道的能量。池底太高的信道不会分配能量。

\begin{example}
\label{ex:ch8-two-parallel-gaussian-e6}
\textbf{两路并联高斯信道，$E=6$}：设两个独立子信道噪声方差为
\[
(\sigma_1^2,\sigma_2^2)=(1,10),
\]
总能量为 $E=6$。求达到容量时的能量分配和容量。

\noindent \textbf{解题流程}：\\[2pt]
步骤一：先假设两个子信道都被使用，检查 $E_i=B-\sigma_i^2$ 是否非负。\\
步骤二：若某个 $E_i<0$，删去对应高噪声子信道。\\
步骤三：重新确定水位并计算容量。

\textbf{解}：若两个子信道都使用，则
\[
(B-1)+(B-10)=6,
\]
得到
\[
B=8.5.
\]
此时
\[
E_2=B-10=-1.5<0,
\]
不符合非负能量分配，所以第二个子信道不用。于是
\[
E_1=6,
\qquad
E_2=0,
\qquad
B=7.
\]
容量为
\[
C=\frac12\log_2\left(1+\frac61\right)
=\frac12\log_2 7
\approx1.404\text{ bit/2自由度}.
\]
\end{example}

\begin{example}
\label{ex:ch8-two-parallel-gaussian-e15}
\textbf{两路并联高斯信道，$E=15$}：噪声方差仍为
\[
(\sigma_1^2,\sigma_2^2)=(1,10),
\]
总能量改为 $E=15$。求达到容量时的能量分配和容量。

\noindent \textbf{解题流程}：\\[2pt]
步骤一：假设两个子信道都使用。\\
步骤二：由总能量求水位 $B$。\\
步骤三：检查能量为正后代入并联容量公式。

\textbf{解}：若两个子信道都使用，则
\[
(B-1)+(B-10)=15,
\]
所以
\[
B=13.
\]
于是
\[
E_1=B-1=12,
\qquad
E_2=B-10=3,
\]
两者均为正。容量为
\[
C=\frac12\log_2\left(1+\frac{12}{1}\right)
+\frac12\log_2\left(1+\frac{3}{10}\right)
\]
\[
=\frac12\log_2 13+\frac12\log_2 1.3
=\frac12\log_2 16.9
\approx2.040\text{ bit/2自由度}.
\]
\end{example}

\subsection{波形信道模型与自由度}

\subsubsection{加性高斯噪声波形信道}

波形信道处理的是连续时间信号。设实信号 $x(t)$ 是输入，线性时不变信道冲激响应为 $g(t)$，频率响应为 $G(f)$，加性高斯噪声为 $z(t)$，输出为
\[
y(t)=x(t)*g(t)+z(t).
\]
若输入信号平均功率受限，则
\[
\frac1T\int_{-T/2}^{T/2}x^2(t)\,dt\le P.
\]

当在信号频带 $B$ 内 $G(f)=1$，且噪声在该频带内为常数谱密度时，信道可化为理想限带 AWGN 信道：
\[
y(t)=x(t)+n(t).
\]
若噪声单边功率谱密度为 $N_0$，带宽为
\[
W=\int_B df,
\]
则频带内噪声功率为 $N_0W$，信噪比为
\[
\frac{P}{N_0W}.
\]

\subsubsection{波形信道的离散化}

波形信道容量的计算依靠等价离散化。对于限时 $T$、限带 $W$ 的实信号，抽样定理说明抽样间隔可取
\[
\frac{1}{2W},
\]
总抽样点数为
\[
N=2TW.
\]
因此限时限带波形信道等价为 $N=2TW$ 个实自由度组成的并联离散时间信道。

这个结论也可以从傅氏级数展开理解：频率间隔为 $1/T$，正负频率一起计算，在带宽 $W$ 内共有约 $2TW$ 个独立系数。无论用时域抽样还是正交展开，得到的自由度数都是
\[
N=2TW.
\]

在时间 $T$ 内，波形信道互信息可以由等价矢量信道互信息表示为
\[
I_T[x(t);y(t)]=\lim_{N\to\infty}I(X^N;Y^N).
\]
平均功率约束变成能量约束
\[
\sum_{i=1}^{N}x_i^2\le PT.
\]
于是波形信道单位时间容量定义为
\[
C=\lim_{T\to\infty}\frac1T\max_{p(x^N):\,\sum_i x_i^2\le PT} I(X^N;Y^N).
\]

\subsection{AWGN 波形信道容量}

限带 AWGN 信道是通信系统中最常见的模型。取抽样率 $2W$ 后，带宽为 $W$ 的白噪声抽样值彼此不相关；由于噪声是高斯的，不相关也意味着独立。这样波形信道就变成 $N=2TW$ 个独立高斯子信道。

每个子信道的噪声方差为
\[
\frac{N_0}{2},
\]
由于所有子信道噪声方差相同，总功率在各自由度上平均分配达到容量。由离散时间 AWGN 容量公式得到单位时间容量。

\begin{formulaBox}{Shannon 限带 AWGN 容量公式}
若 AWGN 信道带宽为 $W$ Hz，输入平均功率为 $P$ W，白噪声单边功率谱密度为 $N_0$ W/Hz，则信道每秒容量为
\[
C=W\log_2\left(1+\frac{P}{N_0W}\right)
\quad\text{bit/s}.
\]
其中 $N_0W$ 是带宽内噪声功率，$P/(N_0W)$ 是接收信噪比。
\end{formulaBox}

也可写成每自由度容量
\[
\frac{C}{2W}=\frac12\log_2\left(1+\frac{P}{N_0W}\right)
\quad\text{bit/自由度},
\]
或频谱效率上限
\[
\frac{C}{W}=\log_2\left(1+\frac{P}{N_0W}\right)
\quad\text{bps/Hz}.
\]

这个公式的使用条件很重要：输入仅受平均功率约束；达到容量的输入应是高斯过程；噪声必须是加性、独立、带内白高斯噪声；$W$ 是正频率带宽，频带可以由若干不相邻频段组成；只关心信号带宽内的噪声谱密度，带外噪声不进入公式。

\begin{example}
\label{ex:ch8-awgn-capacity}
\textbf{AWGN 信道容量计算}：一限带 AWGN 信道带宽为 $1\,\mathrm{MHz}$，信号功率为 $10\,\mathrm{W}$，噪声功率谱为
\[
\frac{N_0}{2}=10^{-9}\,\mathrm{W/Hz}.
\]
求信道容量。

\noindent \textbf{解题流程}：\\[2pt]
步骤一：由 $N_0/2$ 求单边谱密度 $N_0$。\\
步骤二：计算带内噪声功率 $N_0W$。\\
步骤三：代入 $C=W\log_2(1+P/(N_0W))$。

\textbf{解}：由题意
\[
N_0=2\times10^{-9}\,\mathrm{W/Hz},
\qquad
W=10^6\,\mathrm{Hz}.
\]
所以
\[
N_0W=2\times10^{-9}\times10^6=2\times10^{-3}\,\mathrm{W}.
\]
代入 Shannon 公式：
\[
C=10^6\log_2\left(1+\frac{10}{2\times10^{-3}}\right)
=10^6\log_2(5001)
\approx1.23\times10^7\,\mathrm{bps}.
\]
\end{example}

\subsection{高斯信道编码定理}

AWGN 容量不只是一个互信息最大值，它也给出可靠传输的分界线。这个结论和第七章离散有噪信道编码定理形式相同，只是这里还带有输入平均功率约束。

\begin{theoremBox}{高斯信道编码定理}
对于输入平均功率受限的加性高斯噪声信道，设信道容量为 $C$。若传输速率满足
\[
R\le C,
\]
则总可以找到一种编码方式，使差错率任意小。若
\[
R>C,
\]
则不存在能使错误概率任意小的编码方式。
\end{theoremBox}

这个定理说明，Shannon 公式给出的不是某个具体调制方式的性能，而是所有编码和调制方式在理想条件下共同面对的上限。实际系统只能接近它，不能在相同模型和相同约束下突破它。

\subsection{功率利用率与频谱利用率}

通信系统常用两个指标描述性能。功率利用率常用每比特能量与噪声单边功率谱密度之比
\[
\frac{E_b}{N_0}
\]
衡量。它越小，表示用更少能量传输一个可靠比特，功率利用率越高。频谱利用率定义为
\[
\frac{R}{W}
\quad\text{bps/Hz},
\]
表示单位带宽承载的信息速率。

根据高斯信道编码定理，可靠传输要求 $R\le C$，因此
\[
R\le W\log_2\left(1+\frac{P}{N_0W}\right).
\]
又因为
\[
P=E_bR,
\]
代入并整理可得
\[
\frac{R}{W}\le\log_2\left(1+\frac{E_b}{N_0}\frac{R}{W}\right),
\]
等价地，可靠通信必须满足
\[
\frac{E_b}{N_0}\ge
\frac{2^{R/W}-1}{R/W}.
\]

这条曲线把平面分为可靠通信可能区域和不可能区域。为了可靠通信，系统的 $E_b/N_0$ 与 $R/W$ 不能落在曲线以下。它还说明一个重要事实：功率利用率和频谱利用率不能同时无限提高。想降低每比特能量，通常要牺牲频谱效率；想提高频谱效率，通常需要更高的每比特能量。

当
\[
\frac{R}{W}\to0
\]
时，右边极限为
\[
\lim_{R/W\to0}\frac{2^{R/W}-1}{R/W}=\ln2\approx0.693.
\]
换算为 dB：
\[
10\log_{10}(\ln2)=-1.59\,\mathrm{dB}.
\]
这就是 AWGN 信道可靠通信所需的最低能量信噪比，也称 Shannon 极限。它对应无限带宽、频谱效率趋近于 $0$ 的理想情况。

\begin{example}
\label{ex:ch8-reliable-transmission-bandwidth}
\textbf{可靠传输与带宽判断}：已知
\[
\frac{E_b}{N_0}=25\,\mathrm{dB},
\]
问带宽分别为 $100\,\mathrm{kHz}$ 和 $10\,\mathrm{kHz}$ 时，能否可靠传输速率为 $1\,\mathrm{Mbps}$ 的数据。

\noindent \textbf{解题流程}：\\[2pt]
步骤一：分别计算频谱效率 $R/W$。\\
步骤二：由 $E_b/N_0\ge(2^{R/W}-1)/(R/W)$ 求可靠通信所需门限。\\
步骤三：把门限换成 dB 后与 $25\,\mathrm{dB}$ 比较。

\textbf{解}：当 $W=100\,\mathrm{kHz}$ 时，
\[
\frac{R}{W}=\frac{10^6}{10^5}=10.
\]
所需门限为
\[
\frac{E_b}{N_0}\ge\frac{2^{10}-1}{10}=102.3.
\]
换算为 dB：
\[
10\log_{10}(102.3)\approx20.09\,\mathrm{dB}.
\]
由于 $25\,\mathrm{dB}>20.09\,\mathrm{dB}$，通过适当编码可以实现可靠传输。

当 $W=10\,\mathrm{kHz}$ 时，
\[
\frac{R}{W}=\frac{10^6}{10^4}=100.
\]
所需门限为
\[
\frac{E_b}{N_0}\ge\frac{2^{100}-1}{100},
\]
远大于 $25\,\mathrm{dB}$。因此无论采用何种编码方式，都不能在该 AWGN 模型和给定能量信噪比下实现可靠传输。
\end{example}

\subsection{题型整理：第八章常见计算}

\textbf{题型快速判断}：看到题目后先判断它是单自由度信道、并联信道，还是波形信道。
\begin{itemize}
  \item 给 $\{X^N,p(y\mid x),Y^N\}$ 或 $p(y^N\mid x^N)$ $\to$ 先判断是否满足无记忆分解；
  \item 给连续信道容量定义 $\to$ 先找输入约束，是平均能量约束还是幅度约束；
  \item 给 $Y=X+Z$ $\to$ 用 $I(X;Y)=h(Y)-h(Z)$；
  \item 给离散时间 AWGN $\to$ 用 $C=\frac12\log_2(1+E/\sigma_z^2)$；
  \item 给多个高斯子信道和总能量 $\to$ 用注水原理，先求水位再检查非负；
  \item 给带宽、功率、噪声谱密度 $\to$ 用 $C=W\log_2(1+P/(N_0W))$；
  \item 给 $E_b/N_0$ 和 $R/W$ $\to$ 用可靠通信门限判断是否可能。
\end{itemize}

\textbf{答题模板}：
\begin{enumerate}
  \item \textbf{连续信道容量题}：先写约束容量 $C(\beta)=\max I(X;Y)$，再把约束具体化为 $E[X^2]\le E$ 或幅度范围。
  \item \textbf{加性噪声题}：写 $p(y\mid x)=p_Z(y-x)$，再用 $h(Y\mid X)=h(Z)$ 和 $I(X;Y)=h(Y)-h(Z)$。
  \item \textbf{AWGN 题}：确认噪声方差或噪声谱密度，分清离散时间公式和波形公式，最后写清单位。
  \item \textbf{注水题}：先假设所有子信道都使用，求 $B$；若某个 $E_i=B-\sigma_i^2<0$，删去该子信道后重算。
  \item \textbf{可靠性判断题}：先算频谱效率 $R/W$，再求所需 $E_b/N_0$ 门限并换算 dB 比较。
\end{enumerate}

\begin{warningBox}{第八章易错点汇总}
\begin{itemize}
  \item 连续信道容量必须带输入约束，不写功率或幅度约束就直接最大化，通常会得到无意义的无限容量。
  \item 离散时间 AWGN 公式 $\frac12\log_2(1+E/\sigma_z^2)$ 的单位是 bit/自由度，波形 AWGN 公式 $W\log_2(1+P/(N_0W))$ 的单位是 bit/s。
  \item 题目若给 $N_0/2$，代入 Shannon 公式前要先得到单边谱密度 $N_0$。
  \item 注水分配中 $E_i=B-\sigma_i^2$ 只对正能量子信道成立，算出负数时不能保留，必须把该子信道能量置零后重算。
  \item Shannon 极限 $\ln2=-1.59\,\mathrm{dB}$ 是 $R/W\to0$ 时的最低 $E_b/N_0$，不是任意频谱效率下都能达到的门限。
\end{itemize}
\end{warningBox}

\subsection{公式速查表}

{\scriptsize
\noindent\begin{tabularx}{\textwidth}{@{}l>{\raggedright\arraybackslash}X@{}}
\toprule
\textbf{类别} & \textbf{公式 / 核心结论} \\
\midrule
时间离散连续信道 & $\{X^N,p(y\mid x),Y^N\}$，其中 $p(y\mid x)=p(y_1,\dots,y_N\mid x_1,\dots,x_N)$ \\
\midrule
无记忆分解 & \textcolor{Red}{$p(y\mid x)=\prod_{n=1}^{N}p(y_n\mid x_n)$} \\
\midrule
约束容量 & \textcolor{Red}{$C(\beta)=\max_{p(x):\,E[f(X)]\le\beta}I(X;Y)$} \\
\midrule
平均能量约束 & $f(x)=x^2$，$E[X^2]\le E$，$C=\max_{p(x):\,E[X^2]\le E}I(X;Y)$ \\
\midrule
加性噪声信道 & \textcolor{Red}{$Y=X+Z$，$p(y\mid x)=p_Z(y-x)$} \\
\midrule
加性噪声互信息 & \textcolor{Red}{$h(Y\mid X)=h(Z)$，$I(X;Y)=h(Y)-h(Z)$} \\
\midrule
离散时间 AWGN 容量 & \textcolor{Red}{$C=\frac12\log_2\left(1+\frac{E}{\sigma_z^2}\right)$ bit/自由度} \\
\midrule
并联高斯信道 & $C=\sum_i\frac12\log_2\left(1+\frac{E_i}{\sigma_i^2}\right)$，$\sum_iE_i\le E$ \\
\midrule
注水原理 & \textcolor{Red}{$E_i=B-\sigma_i^2$ 若 $B>\sigma_i^2$，否则 $E_i=0$} \\
\midrule
波形自由度 & \textcolor{Red}{$N=2TW$}，限时 $T$、限带 $W$ 的实波形有 $2TW$ 个自由度 \\
\midrule
AWGN Shannon 公式 & \textcolor{Red}{$C=W\log_2\left(1+\frac{P}{N_0W}\right)$ bit/s} \\
\midrule
频谱效率上限 & $C/W=\log_2\left(1+\frac{P}{N_0W}\right)$，单位 bps/Hz \\
\midrule
高斯信道编码定理 & \textcolor{Red}{$R\le C$ 时可使差错率任意小，$R>C$ 时不能可靠传输} \\
\midrule
功率与频谱效率 & \textcolor{Red}{$\dfrac{E_b}{N_0}\ge\dfrac{2^{R/W}-1}{R/W}$} \\
\midrule
Shannon 极限 & \textcolor{Red}{$\lim_{R/W\to0}\dfrac{2^{R/W}-1}{R/W}=\ln2=-1.59\,\mathrm{dB}$} \\
\bottomrule
\end{tabularx}
}

\subsection{本章要求}

本章学习重点：
\begin{itemize}
  \item 理解时间离散连续信道模型 $\{X^N,p(y\mid x),Y^N\}$；
  \item 掌握平稳无记忆连续信道的转移概率密度分解；
  \item 理解连续信道容量必须带输入功率或幅度约束；
  \item 会写约束容量 $C(\beta)=\max_{p(x):E[f(X)]\le\beta}I(X;Y)$；
  \item 掌握平均能量约束 $E[X^2]\le E$ 下的容量表达；
  \item 掌握加性噪声信道 $Y=X+Z$、$p(y\mid x)=p_Z(y-x)$ 和 $I(X;Y)=h(Y)-h(Z)$；
  \item 会计算幅度受限均匀噪声例题的最大互信息；
  \item 掌握离散时间 AWGN 容量 $C=\frac12\log_2(1+E/\sigma_z^2)$；
  \item 理解并联高斯信道的注水原理，并会处理两路信道能量分配；
  \item 理解波形信道离散化和 $N=2TW$ 个自由度的来源；
  \item 掌握 Shannon 公式 $C=W\log_2(1+P/(N_0W))$ 的单位和使用条件；
  \item 理解高斯信道编码定理给出的可靠传输分界；
  \item 掌握功率利用率和频谱利用率关系式；
  \item 记住 Shannon 极限 $\ln2=-1.59\,\mathrm{dB}$ 的含义和适用条件。
\end{itemize}
